\documentclass[UTF8]{ctexart}
\usepackage{xeCJK}
\usepackage{geometry}
\geometry{left=2cm,right=2cm,top=2.5cm,bottom=2.5cm}
\setlength{\headheight}{12.7pt}

\usepackage{titlesec}
\titleformat{\section}
  {\normalfont\large\bfseries}
  {\thesection}
  {1em}
  {}
\titlespacing*{\section}
  {0pt}
  {3.5ex plus 1ex minus .2ex}
  {2.3ex plus .2ex}

\usepackage{amsmath}
\usepackage{amssymb}
\usepackage{amsthm}
\usepackage{enumitem}
\setlist[enumerate]{itemsep=0pt,topsep=0pt,parsep=0pt,leftmargin=0pt,itemindent=2.5em}
\setlist[itemize]{itemsep=0pt,topsep=0pt,parsep=0pt,leftmargin=0pt,itemindent=2.5em}
\usepackage{fancyhdr}
\pagestyle{fancy}
\fancyhf{}
\usepackage{graphicx}
\usepackage{hyperref}
\usepackage{indentfirst}
\usepackage{lmodern}


\begin{document}
\setlength{\abovedisplayskip}{1pt}
\setlength{\belowdisplayskip}{1pt}

\section{量词消去判别法}

\hypertarget{sec:定理1.1}{}
\noindent\textbf{定理1.1}

给定一阶语言$\mathcal{L}$, $\mathcal{L}$-理论$T$和公式$\varphi$, 以下等价:
\begin{enumerate}
    \item 对$T$的任意模型$A,B$, 如果有它们的子结构$M,N$以及同构$p:M\cong N$,
    那么
    \[
        \forall\sigma:\mathcal{V}\to M:\quad
        (A,\sigma)\models\varphi\leftrightarrow(B,p\circ\sigma)\models\varphi.
        \\[3pt]
    \]
    \item 存在无量词公式$\psi$使得$\mathrm{FV}(\psi)\subset\mathrm{FV}(\varphi)$
    并且$\varphi\overset{T}{\leftrightarrow}{\psi}$.
\end{enumerate}

\noindent{\kaishu 证明.}

\noindent{\kaishu 1 $\to$ 2.}

假设1.条件成立. 准备工作: 记$W=\mathrm{FV}(\varphi)$, 以及:
\[\begin{aligned}
    &Q&&=\{\psi\in F^{\mathcal{L}}\mid
    \psi\text{\,无量词}\land\mathrm{FV}(\psi)\subset W\},\\[-3pt]
    &P&&=\{\psi\in Q\mid T\cup\{\varphi\}\models\psi\}.
\end{aligned}\]

首先证明$T\cup P\models\varphi$, 即对任意的$(A,\sigma)\models T\cup P$,
下证$(A,\sigma)\models\varphi$. 这时再记
\[
    P_A=\{\psi\in Q\mid(A,\sigma)\models\psi\},
\]
因为$(A,\sigma)\models T\cup P$, 所以显然有$P\subset P_A$.

第一, 证明$T\cup\{\varphi\}\cup P_A$一致.
任取有限多个$\psi_1,\cdots,\psi_n\in P_A$, 易证也有$\displaystyle\psi=
\bigwedge_{i=1}^n\psi_i\in P_A$.

\vspace{3pt}
假设$T\cup\{\varphi,\psi\}$不一致, 则$T\cup\{\varphi\}\models\neg\psi$,
依定义$\neg\psi\in P\subset P_A$, 与$\psi\in P_A$矛盾.
所以$T\cup\{\varphi,\psi\}$一致, 即$T\cup\{\varphi,\psi_1,\cdots,\psi_n\}$一致.
根据\href{../03 一阶紧致性定理/03 一阶紧致性定理#nameddest=sec:定理1.1}{紧致性定理},
$T\cup\{\varphi\}\cup P_A$一致, 取$(B,\tau)\models T\cup\{\varphi\}\cup P_A$.

第二, 对任意满足$\mathrm{FV}(\psi)\subset W$的原子公式$\psi$,
显然$\psi\in Q$, 同时$\neg\psi\in Q$. 此时:
\begin{enumerate}
    \item 如果$(A,\sigma)\models\psi$, 那么$\psi\in P_A$, 于是$(B,\tau)\models\psi$.
    \item 如果$(A,\sigma)\not\models\psi$, 那么$\neg\psi\in P_A$,
    于是$(B,\tau)\models\neg\psi$, 即$(B,\tau)\not\models\psi$.
\end{enumerate}
综上, 应用\href{../01 一阶语言的结构/04 生成子结构#nameddest=sec:定理1.4}{引理},
记$M=\langle\,\sigma[W]\,\rangle,N=\langle\,\tau[W]\,\rangle$, 则有同构
$p:M\cong N$使得$p\circ\sigma\!\restriction_W=\tau\!\restriction_W$.

第三, 取一个$\sigma':\mathcal{V}\to M$使得$\sigma'\!\restriction_W=\sigma\!\restriction_W$,
令$\tau'=p\circ\sigma'$, 则显然也有$\tau'\!\restriction_W=\tau\!\restriction_W$. 再根据1.条件有
\[
    (B,\tau)\models\varphi\implies(B,\tau')\models\varphi\implies
    (A,\sigma')\models\varphi\implies(A,\sigma)\models\varphi.
\]
至此已证$(A,\sigma)\models\varphi$, 即已证$T\cup P\models\varphi$.

然后再根据紧致性定理, 存在有限多个$\psi_1,\cdots,\psi_n\in P$使得
$T\cup\{\psi_1,\cdots,\psi_n\}\models\varphi$.
令$\displaystyle\psi=\bigwedge_{i=1}^n\psi_i$, 则:
\begin{enumerate}
    \item $\psi$也是无量词公式, 并且$\mathrm{FV}(\psi)\subset W$.
    \item 对每个$\psi_i$有$T\cup\{\varphi\}\models\psi_i$,
    从而$T\cup\{\varphi\}\models\psi$.
    \item $T\cup\{\psi\}\models T\cup\{\psi_1,\cdots,\psi_n\}\models\varphi$.
\end{enumerate}
综上, $\psi$满足2.条件.

\noindent{\kaishu 2 $\to$ 1.}

假设存在无量词公式$\psi$使得$\varphi\overset{T}{\leftrightarrow}\psi$.

对$T$的任意模型$A,B$, 如果有它们的子结构$M,N$以及同构$p:M\cong N$,
则对任意的$\sigma:\mathcal{V}\to M$有
\begin{equation*}\begin{aligned}
    &(A,\sigma)\models\varphi\iff(A,\sigma)\models\psi\iff
    (M,\sigma)\models\psi\\[-3pt]
    \iff\,&(N,p\circ\sigma)\models\psi\iff(B,p\circ\sigma)\models\psi\iff
    (B,p\circ\sigma)\models\varphi.
\end{aligned}\tag*{\qedsymbol}\end{equation*}





\newpage

\hypertarget{sec:定理1.2}{}
\noindent\textbf{定理1.2 (量词消去判别法)}

给定一阶语言$\mathcal{L}$和$\mathcal{L}$-理论$T$,
$T$是\href{01 量词消去理论#nameddest=sec:定义1.2}{量词消去理论}当且仅当:

对任意的无量词公式$\varphi$、变元$x$、$T$的任意模型$A,B$,
如果有它们的子结构$M,N$以及同构$p:M\cong N$, 那么
\[
    \forall\sigma:\mathcal{V}\to M:\quad
    (A,\sigma)\models\exists x\varphi\leftrightarrow
    (B,p\circ\sigma)\models\exists x\varphi.
\]

\noindent{\kaishu 证明.}

由\href{01 量词消去理论#nameddest=sec:定理1.2}{引理}和上述定理即得.
\hfill $\qedsymbol$

% 模型论的学习暂告一段落, 据说只算是入门水平. 谢谢黄书棋老师, 谢谢 bilibili.
% 之后计划尽快地回到量子信息理论相关的内容. 证明论也在计划内.
% 2026.03.17

\end{document}